\documentclass[12pt]{article}
\usepackage[utf8]{inputenc}
\usepackage[T1]{fontenc}
\usepackage{hyperref}
\title{Papantla Totonac Language Resources for K-12 Curriculum Development}
\author{Compiled by Claude Sonnet 5 (claude-sonnet-5) for Andr\'as Kornai}
\date{}
\begin{document}
\maketitle

\section{Introduction}
Papantla Totonac (Tutunak\'u; ISO 639-3 \texttt{top}; Glottocode papa1238) is a Totonacan language of Veracruz, Mexico. Sources in the search gave around 80,000 speakers, and the reference description gives about 100,000. The search found no complete, current, online dictionary or grammar for this variety.

\section{Primary Educational Resources}
No primers, textbooks or school programmes for Papantla Totonac appeared in the search results.

\section{Linguistic and Reference Materials}
\begin{itemize}
\item Bibliography of Totonac linguistic studies (Sam Holley-Kline). It lists Papantla works including Levy, \emph{Totonaco de Papantla, Veracruz} (1990), Aschmann, \emph{Diccionario Totonaco de Papantla, Veracruz} (1973), and Garc\'ia Ramos, \emph{Manual de fon\'etica totonaca} (2006). I did not verify that any of these can be accessed. \url{https://www.bibliografiatotonaca.com/estudios-linguumliacutesticos.html}
\item Omniglot page on the Papantla Totonac writing system. \url{https://omniglot.com/writing/papantlatotonac.htm}
\item \emph{Upper Necaxa Totonac Dictionary} (Beck, 2011). This covers a different Totonac variety and is listed here only as a related lexical resource. \url{https://www.degruyterbrill.com/document/doi/10.1515/9783110238235/html}
\end{itemize}

\section{Religious and Community Resources}
\begin{itemize}
\item Endangered Languages Alliance page on Totonac (overview). \url{https://www.elalliance.org/languages/totonac}
\end{itemize}

\section{Digital and Audio Resources}
No audio or software for Papantla Totonac turned up in this search.

\section{Curriculum Development Context}
The language is thought to be vigorous, but no school materials were found. A curriculum would have to start from the older SIL and El Colegio de M\'exico descriptions, the Omniglot orthography notes and work with the community. A more targeted search of Mexican education-ministry and INALI materials is needed.

\end{document}